\section{Empirical analysis}
\label{sec:main_results}
\noindent In this section, we present our empirical strategy and results. We begin by outlining the debt overhang framework from \cite{Hennessy2004}, which guides our empirical design and generates predictions for both investment and payout policy. We then investigate how uncertainty about gold clause enforcement---and the resulting uncertainty about firms' effective leverage---affected corporate decisions in 1933 and 1934.

\subsection{Conceptual framework}
\label{sec:concept}
\noindent To motivate our empirical strategy, let's begin with the standard investment regression relating investment to average $Q$:
\begin{align}
\label{eq:classic_inv_reg}
i_{j,t} = \beta Q_{j,t-1} + \epsilon_{j,t},
\end{align}
where $i_{j,t} = I_{j,t}/K_{j,t-1}$ is the investment rate for firm $j$ in year $t$, and $Q_{j,t} = V_{j,t}/K_{j,t}$ is average $Q$, with $I$ denoting investment, $V$ firm value, and $K$ the replacement cost of capital. The theoretical foundation for this specification comes from \cite{Hayashi1982} and \cite{Abel1994}, who show that, under certain conditions, average $Q$ is a sufficient statistic for investment policy.

For firms with long-term debt, however, \cite{Hennessy2004} shows that capital structure creates a wedge $r_{j,t}$ between a firm's \textit{average} $Q$ and its \textit{marginal} $q$, which governs the investment policy that maximizes the value of equityholders:
\begin{align}
\label{eq:overhang_wedge}
q_{j,t} = Q_{j,t} - r_{j,t}.
\end{align}
The wedge is given by $r_{j,t} = \frac{R_{j,t}}{K_{j,t}}$, with $R_{j,t}$ denoting the expected present value of debtholder's recovery upon default. This wedge arises because recovery value, although part of the firm's total value, does not accrue to equityholders, who therefore ignore its contribution to average $Q$ when making investment decisions. To account for this effect, \cite{Hennessy2004} shows that the standard investment regression must be modified as follows:
\begin{align}
\label{eq:overhang_inv_reg}
i_{j,t} = \beta Q_{j,t-1} - \beta r_{j,t-1} + \epsilon_{j,t}.
\end{align}

The gold clause litigation provides a natural setting to apply this framework. A ruling against the abrogation would have raised the claim of debtholders on the firm's assets and increased the probability of default. Accordingly, the legal uncertainty about the abrogation must have increased the expected recovery value $R_{j,t}$ and thus the debt overhang wedge $r_{j,t}$. Crucially, this effect varied across firms in proportion to their exposure to gold-denominated debt, thereby generating differential increases in debt overhang across firms during the litigation period.

Empirically, measuring the wedge $r_{j,t}$ directly is challenging since it depends on default probabilities and expected recovery rates, which are not directly observable. We therefore parameterize the wedge as a time-dependent linear function of our gold exposure measure $\tilde{d}_{j,t}$, which captures the share of a firm's liabilities subject to gold clauses:
\begin{align}
\label{eq:parameterization}
i_{j,t} = \beta Q_{j,t-1} + \sum_{t} \gamma_{t} \tilde{d}_{j,t-1} + f_{j} + \tau_{t} + \epsilon_{j,t},
\end{align}
where $f_{j}$ and $\tau_{t}$ are firm and time fixed-effects, respectively. This specification allows for flexible, time-varying effects of gold exposure on investment, which can be interpreted as a generalized difference-in-differences design around the litigation period with $\tilde{d}_{j,t}$ used as a continuous treatment. Under the hypothesis that gold clause uncertainty increased the debt overhang wedge in 1933 and 1934, we expect the corresponding coefficients, $\gamma_{1933}$ and $\gamma_{1934}$, to be negative. Conversely, the coefficients $\gamma_t$ for years outside the litigation period should be close to zero.

\subsection{Investment}
\label{sec:investment}
\noindent Table \ref{tab:inv_main} reports our net investment regression results. All specifications include firm and year fixed effects, with 1932 (the year immediately preceding the abrogation) as the omitted category. Column 1 presents the estimates of the classical investment regression from equation (\ref{eq:classic_inv_reg}). Consistent with Q-theory and with empirical evidence from more recent samples (e.g., \cite{Peters2017}), average Q is positively associated with investment.

Column 2 presents our main debt overhang results from equation (\ref{eq:overhang_inv_reg}). The $\text{year}\times\tilde{d}$ interaction coefficients are significantly negative only in 1933 and 1934---the years when the gold clause was litigated---while remaining statistically insignificant in all other years. The economic magnitudes are substantial: fully exposed firms ($\tilde{d}=1$) reduced investment by an additional 6.0 percentage points in 1933 and 4.1 percentage points in 1934 relative to unexposed firms ($\tilde{d}=0$). The interaction coefficient in 1935 is marginally significant and positive, suggesting that firms with gold exposure partially caught up with unexposed firms after the Supreme Court decision. Notably, the investment decline we document is concentrated in 1933, with a smaller---though still statistically significant---effect in 1934. This pattern is consistent with gradual resolution of uncertainty during the litigation period: as lower courts ruled in favor of abrogation throughout 1933 and early 1934, firms may have updated their expectations about the likelihood that the Supreme Court would ultimately uphold the gold clause.

To visually examine pre-trends, Figure \ref{fig:gold_coeffs} plots the yearly interaction coefficients reported in column 2 of Table \ref{tab:inv_main}. Before 1933, the interaction coefficients are not statistically significant and show no visible trend. Moreover, following the Supreme Court decision to uphold the abrogation in 1935, the effect of $\tilde{d}$ on investment quickly dissipates. This pattern---no pre-trends through 1932, sharp effects during the 1933-1934 litigation period, and subsequent reversal after the 1935 Supreme Court ruling---suggests that gold exposure depressed investment during the window of legal uncertainty and ceased to have an effect once the constitutional question was resolved.

The absence of effects in other years is particularly noteworthy for 1937 and 1938, which were also years of economic distress. If gold exposure $\tilde{d}$ captured inherent firm characteristics that make some companies more vulnerable during downturns, we would expect to observe differential investment patterns whenever economic conditions deteriorate. The fact that we find no effect of $\tilde{d}$ on investment during this subsequent recession, three years after the Supreme Court's ruling, suggests that our 1933--1934 findings reflect the specific impact of gold clause litigation rather than persistent differences in firms' sensitivity to adverse economic conditions.

While the evidence thus far supports a debt overhang interpretation, we next examine whether alternative mechanisms could explain the differential investment patterns between exposed and unexposed firms. First, we consider the possibility that gold exposure instead captures firms' financial constraints related to the contraction of credit supply. This concern is motivated by Figure \ref{fig:bond_issuance}, which shows that during the litigation of the gold clause, the issuance of corporate bonds was virtually inexistent. If this drastic reduction in primary market activity is attributable to a contraction in credit supply, firms with bonds maturing during these years may have faced severe difficulties rolling over their debt and thus reduced investment. Indeed, \cite{Benmelech2018} show that disruptions in intermediation and bond market dislocations were key drivers of the investment contraction in the early years of the Great Depression, even for large firms. 

To test if this rollover channel is driving our results, we re-estimate equation (\ref{eq:overhang_inv_reg}) excluding firms with bonds maturing at any point between 1931 and 1934, inclusively. Column 3 of Table \ref{tab:inv_main} reports results consistent with the baseline specification: the interaction coefficients for 1933 and 1934 are -0.058 and -0.037, respectively, compared to -0.060 and -0.041 in the baseline specification. The stability of the coefficient estimates shows that the effect of gold exposure is present among firms with no immediate refinancing needs, suggesting that $\tilde{d}$ is not simply capturing firms' vulnerability to credit constraints.

We next consider whether gold-exposed firms reduced investment because they redirected resources toward redeeming gold-denominated debt during the litigation period. This interpretation is plausible given the decline in the number of gold bonds from 250 in 1930 to 212 in 1933 and 189 in 1934, as documented in Table \ref{tab:bond_stats}. To test whether this channel is driving our results, we re-estimate equation (\ref{eq:overhang_inv_reg}) excluding firms that repurchased their bond issues in 1933 and 1934. Column 4 reports interaction coefficients for 1933 and 1934 of -0.068 and -0.032, respectively, which remain economically and statistically significant. These results indicate that the decline in investment among gold-exposed firms extends beyond those actively managing their debt obligations through redemptions during the litigation period.

Relatedly, we consider whether our results might be driven by firms without long-term liabilities, which could differ systematically from firms with long-term obligations on time-varying unobservable characteristics. Column 5 of Table \ref{tab:inv_main} reports results from re-estimating the baseline overhang specification restricting the sample to firms with a positive amount of long-term liabilities. The interaction coefficients for 1933 and 1934 are $-0.058$ and $-0.043$, respectively, which are statistically significant and remain consistent with the negative coefficients observed in the baseline specification. These results indicate that the baseline findings are robust to excluding firms with zero long-term liabilities, alleviating concerns that our results are driven by systematic differences between firms with and without long-term obligations.

Having established that our findings are robust to various sample restrictions, we next conduct placebo tests to verify that the documented investment effect is specifically attributable to gold clause exposure. To do so, we re-estimate equation (\ref{eq:overhang_inv_reg}) while replacing the numerator of $\tilde{d}$ with the book value of preferred shares. If the investment effect documented in our baseline specification results from gold exposure, we do not expect to find any effect when replacing the numerator of $\tilde{d}$ with preferred shares since the gold clause was not included in those securities. The results in column 6 support this conjecture: the interaction coefficients in 1933 and 1934 are positive and not statistically significant. As a second placebo test, we replace the numerator of $\tilde{d}$ with bank debt. Here again, the results in column 7 show that the 1933 and 1934 interaction terms are positive and not statistically significant.

\subsection{Payouts and other outcomes}
\label{sec:payouts}
\noindent
According to the dynamic debt overhang theory of \cite{Hennessy2004}, the shareholders of distressed firms may optimally increase the equity payouts to extract value from the firm in anticipation of losing their claim on the assets.\footnote{See Section I.C in \cite{Hennessy2004}.} This prediction distinguishes debt overhang from theories of financial frictions, which also predict reduced investment during distress but imply \textit{lower} payouts.

To test this prediction, we first calculate the total equity payout of a firm $j$ at time $t$ as\footnote{We fix the denominator at its 1930 value (or the first available year after 1930) to prevent post-treatment changes in equity structure from distorting the payout measure. Since the gold clause shock may itself affect equity issuance and repurchase decisions, using a time-varying denominator could conflate changes in actual payouts with changes in the capital structure.}
\begin{align*}
    \text{Payout}_{j,t} = \frac{\text{Dividend}_{j,t} + \text{Net repurchase}_{j,t}}{\text{Book value of common stocks}_{j,1930}}.
\end{align*}
We then re-estimate equation (\ref{eq:overhang_inv_reg}) using payouts as the dependent variable. Column 1 of Table \ref{tab:other_outcomes} shows that the interaction coefficients for 1933 and 1934 are 0.038 and 0.054, both statistically significant at the 5\% and 1\% levels, respectively. Outside the litigation period, the coefficients are close to zero and statistically insignificant. Thus, consistent with debt overhang, the degree of gold exposure is associated with greater payouts during the gold clause litigation period, and during this period only.

Columns 2 and 3 investigate which component of payouts drives the results. The dividend regression (column 2) shows statistically significant interaction coefficients of 0.028 and 0.036 for 1933 and 1934, respectively. In contrast, the net repurchase regression (column 3) yields statistically insignificant coefficients for the corresponding years. This pattern in the payout effect likely reflects that, before 1982, dividend payments were the primary driver of payout policy \citep[e.g.,][]{Michaely2002, Skinner2008}.

So far, our payout results are consistent with the hypothesis that the lower investment rates of gold-exposed firms in 1933 and 1934 are driven by debt overhang. However, it could also be that gold-exposed firms increased payouts in 1933 and 1934 simply because these firms were more profitable in those years. To address this concern, we re-estimate equation (\ref{eq:overhang_inv_reg}) using profitability as the dependent variable. The estimation results reported in column 4 show no evidence that gold exposure is associated with increased profitability in 1933 or 1934. If the higher payouts instead reflected greater available resources rather than debt overhang incentives, we would also expect gold-exposed firms to have higher cash holdings or lower leverage during the litigation period. Columns 5 and 6 show that neither cash nor leverage is systematically associated with gold exposure in 1933 or 1934. Together, these results suggest that the payout effect is driven by debt overhang rather than by profitability or resource availability.

In the Internet Appendix, we probe the dividend results along two additional dimensions. First, Table \ref{tabapp:divadd} examines the robustness of the annual dividend findings to sample restrictions and alternative normalizations. The dividend increase during the litigation period is concentrated among firms that were already paying dividends: columns 1 and 3 show strong and statistically significant effects for prior dividend payers, while columns 2 and 4 show no effect for non-payers. This pattern is consistent with the well-documented stickiness of dividend policy. The effect also remains significant across several alternative normalizations of the dependent variable, including dividends scaled by book equity (column 6) and by market capitalization (column 7).

Second, we examine whether the dividend effect holds at the quarterly frequency (Table \ref{tabapp:quarterly_div}). Because dividends exhibit strong seasonal patterns and are sticky, we estimate our baseline specification separately for each calendar quarter, using quarterly cash dividends as the dependent variable. Although the statistical evidence is somewhat weaker at this higher frequency, the quarterly results are concentrated in the first and second quarters. The Q2 interaction coefficients for 1933 and 1934 are both 0.010, statistically significant at the 1\% and 5\% levels, respectively.

In sum, gold-exposed firms in 1933 and 1934 reduced investment and increased payouts---primarily through cash dividends---without changes in profitability, cash, or leverage that would indicate differential financial constraints or greater resource availability.

\subsection{Credit ratings}
\label{sec:rating}

\noindent Standard theories of debt overhang such as \cite{Myers1977} and \cite{Hennessy2004} predict that investment distortions should be more severe for firms closer to default. To test this hypothesis, we collect bond ratings from the Moody's Manuals and construct the dummy variable, ``Low'', which is equal to one if a firm's bond rating is Ba or below in 1930.\footnote{We chose Ba as cutoff since it is the median credit rating of the bond issuers in our sample as of 1930. We use 1930 ratings---rather than contemporaneous ratings---to avoid endogeneity concerns.} We then modify our investment regression Equation~(\ref{eq:parameterization}) to include an interaction term between gold exposure and low credit ratings:
\begin{align}
\label{eq:parameterization_by_ratings}
i_{j,t} = \beta Q_{j,t-1} + \sum_{t} \gamma_{t} \tilde{d}_{j,t-1} + \sum_{t} \theta_{t} \mathbbm{1}_{\text{Low}_{j}} + \sum_{t} \kappa_{t} \tilde{d}_{j,t-1} \times \mathbbm{1}_{\text{Low}_{j}} + f_{j} + \tau_{t} + \epsilon_{j,t},
\end{align}

To ease the interpretation of the coefficients, we normalize $\tilde{d}$ by subtracting the median value of $\tilde{d}$ among firms with positive $\tilde{d}$, which is 0.60. This normalization centers the gold exposure variable at a meaningful reference point---the median exposure level among firms with gold-denominated debt. With this normalization, $\theta_t$ now represents the difference in investment between low-rated and high-rated firms with median exposure to gold. The interpretation of the coefficients $\gamma_t$ and $\kappa_t$ is unaffected by the normalization.

The first column of Table \ref{tab:credit_rating} reports our estimation results of equation (\ref{eq:parameterization_by_ratings}).\footnote{To conserve space, the table reports only the 1933 and 1934 interaction coefficients. The full set of year-by-year results is available in Internet Appendix Table \ref{tabapp:credit_rating}.} The coefficients on $1933 \times \tilde{d}$ and $1934 \times \tilde{d}$ are -0.036 and -0.047, respectively, both statistically significant at the 1\% level. For comparison, the corresponding coefficients in the baseline specification from Table \ref{tab:inv_main} are -0.060 and -0.041. These coefficients represent the marginal effect of gold exposure on investment for high-rated firms. These results indicate that debt overhang depressed investment even among high-rated firms, not just those closest to default.

Turning to the coefficients on the interaction between year indicators and low credit rating, recall that $\theta_t$ represents the difference in investment between low-rated and high-rated firms at the median gold exposure level. The coefficient estimates are $\theta_{1933} = -0.038$ (significant at the 5\% level) and $\theta_{1934} = 0.029$ (significant at the 10\% level). The negative coefficient in 1933 indicates that low-rated firms with median gold exposure reduced investment by an additional 3.8 percentage points relative to high-rated firms, consistent with the prediction that debt overhang effects are more severe for firms closer to default. The positive coefficient in 1934, however, is inconsistent with this prediction, though we note that it is only marginally statistically significant.

Finally, the triple interaction coefficient $\kappa_t$ captures the additional marginal effect of gold exposure for low-rated firms relative to high-rated firms---that is, the difference in how investment responds to a one-unit increase in gold exposure across credit rating groups. The coefficient estimates are $\kappa_{1933} = -0.073$ and $\kappa_{1934} = -0.077$, neither statistically significant. The negative signs are consistent with the debt overhang prediction that low-rated firms should exhibit a stronger marginal response to gold exposure, and the economic magnitudes are substantial. However, the coefficients are imprecisely estimated, so we interpret this evidence with caution.

The second column of Table \ref{tab:credit_rating} reports our estimation results of equation (\ref{eq:parameterization_by_ratings}) with dividends as the dependent variable. The coefficients on $1933 \times \tilde{d}$ and $1934 \times \tilde{d}$ are 0.025 and 0.035, respectively, both statistically significant, compared to 0.028 and 0.036 in the baseline specification from Table \ref{tab:other_outcomes}. The positive and significant estimates indicate that, even among high rated firms, debt overhang induced dividend increases consistent with debt overhang.

As before, the coefficient $\theta_t$ captures the dividend differential between low-rated and high-rated firms at median gold exposure. The coefficient estimates are $\theta_{1933} = 0.013$ (significant at the 5\% level) and $\theta_{1934} = 0.011$ (significant at the 10\% level). The positive coefficients indicate that, relative to high-rated firms, low-rated firms with median gold exposure increased dividends by an additional 1.3 and 1.1 percentage points in 1933 and 1934, respectively. This pattern is consistent with the prediction that debt overhang effects are more severe for firms closer to default, as these firms face stronger incentives to extract value through payouts.

Finally, the triple interaction coefficient $\kappa_t$ captures the additional marginal effect of gold exposure for low-rated firms relative to high-rated firms. The coefficient estimates are $\kappa_{1933} = -0.025$ (not statistically significant) and $\kappa_{1934} = -0.037$ (significant at the 5\% level). The negative signs suggest that the marginal effect of gold exposure on dividends diminishes with the degree of exposure, which is inconsistent with the prediction that low-rated firms should exhibit a stronger response for a given degree of exposure. However, this pattern may reflect non-linearity in the relationship between gold exposure and default risk. At moderate exposure levels, low-rated firms may be the only ones facing sufficient default risk to increase dividends aggressively, whereas at high exposure levels, both low- and high-rated firms face substantial default risk, narrowing the gap between them. Given the limited statistical significance of the triple interaction coefficients, however, we interpret this evidence with caution.

In sum, the credit rating results provide mixed but generally supportive evidence for the hypothesis that debt overhang effects are more pronounced for firms closer to default. Low-rated firms reduced investment more in 1933 and increased dividends more in both 1933 and 1934, consistent with stronger debt overhang effects for firms with weaker credit profiles. However, the triple interaction coefficients capturing the differential \textit{marginal} effect of gold exposure for low-rated firms are either imprecisely estimated (for investment) or suggest a pattern that may reflect non-linearities rather than a simple differential response (for dividends). We therefore interpret these findings as suggestive rather than definitive evidence on the role of credit ratings in mediating debt overhang effects.

\subsection{Controls and robustness}
\label{sec:cont}
\noindent
Although the evidence presented thus far is broadly consistent with the hypothesis that gold exposure depressed firm investment through a debt overhang channel, our study period coincides with major economic and policy changes, including the New Deal and the Agricultural Adjustment Act. These policy shocks pose empirical challenges because they may confound the estimated effect of gold exposure. We therefore conduct several robustness checks to ensure that our documented effects are driven by debt overhang rather than these concurrent shocks.

We first control for the impact of New Deal policies by including industry-year fixed effects in our specification. Major interventions such as the National Industrial Recovery Act's industry codes and the Agricultural Adjustment Act operated primarily at the industry level (e.g., \cite{Cole2004}). By absorbing time-varying, industry-specific shocks, industry-year fixed effects ensure that the gold exposure coefficient is identified solely from within-industry variation, netting out any industry-level policy effects. Given our limited sample size, we group years into four periods to maintain statistical power: the pre-treatment period (1926--1932, omitted), 1933, 1934, and the recovery years (1935--1940). We then interact these periods with 2-digit Standard Industrial Classification (SIC) codes.

The first column of Table \ref{tab:controls} reports the estimation results of the specification with industry-year fixed effects. The 1933 and 1934 interaction coefficients are -0.059 and -0.036, respectively, compared to -0.060 and -0.041 in the baseline. Both coefficients remain statistically significant. The stability of these estimates after absorbing industry-specific time shocks suggests that our results are unlikely to be driven by New Deal policies or other industry-specific disturbances. These results strengthen our interpretation that gold exposure depressed investment through a firm-specific debt overhang channel rather than through industry-wide regulatory or policy changes.

While industry-year fixed effects address industry-level confounds, they do not account for firm-level heterogeneity that may correlate with both gold exposure and investment. To address this concern, we augment equation (\ref{eq:parameterization}) with interaction terms between pre-treatment firm characteristics and time fixed effects. This approach allows firms with different observable characteristics to follow different investment trajectories, isolating the effect of gold exposure from policies that may have differentially affected firms based on these characteristics. For example, if New Deal policies disproportionately affected large firms, highly leveraged firms, or firms with certain financial profiles, allowing these characteristics to have time-varying effects ensures that our gold exposure estimates are not confounded by such differential policy impacts.

Column 2 of Table \ref{tab:controls} reports results from augmenting equation (\ref{eq:parameterization}) with interactions between time indicators and linear controls for pre-treatment firm characteristics: log assets, Tobin's Q, net income over assets, cash holdings, payout ratios, book leverage, market leverage, and log long-term liabilities. This specification uses year fixed effects rather than industry-year fixed effects because the characteristic-period interactions and industry-year cells absorb overlapping variation; including both simultaneously causes several interactions to be dropped for multicollinearity, undermining the purpose of the exercise. The coefficient on the 1933-gold exposure interaction remains negative and statistically significant at -0.035, while the 1934 coefficient is -0.026, maintaining the expected negative sign though losing statistical significance. The loss in significance for 1934 is not surprising given the demanding nature of this specification, which allows eight different firm characteristics to have time-varying effects.

To ensure our results are not sensitive to the functional form of the controls, we next employ a more flexible approach by controlling for deciles of each firm characteristic interacted with time fixed effects. This non-parametric specification allows for non-linear relationships between firm characteristics and investment trajectories during the Depression. Rather than including all characteristics simultaneously---which would require estimating an impractical number of parameters given our sample size---we run eight separate regressions, each controlling for firm and year fixed effects and deciles of a single characteristic interacted with time. Columns 3--10 of Table \ref{tab:controls} report these results.

The findings are consistent with our earlier specifications: all eight estimated coefficients on the 1933-gold exposure interaction are negative and statistically significant. The 1934 coefficients are likewise all negative, with four of the eight specifications retaining statistical significance. This pattern reinforces our main findings while demonstrating that the debt overhang effect, particularly in 1933, is robust to flexible controls for observable firm heterogeneity. In Internet Appendix Table \ref{tabapp:retcontrol}, we also verify robustness to stock market controls based on 1930 returns (average, volatility, and Fama-French betas). Internet Appendix Table \ref{tabapp:controls_indyear} re-estimates columns 2--10 replacing year fixed effects with industry-year fixed effects. The average exposure effect $\tilde{d}$ remains negative and significant at the 5\% level in all nine specifications, and $1933 \times \tilde{d}$ is significant in eight of nine. $1934 \times \tilde{d}$ loses precision in six of nine specifications, as the portfolio decile controls and industry-year cells absorb overlapping variation when included simultaneously.

\subsection{Aggregation}
\label{sec:agg}
\noindent As a final exercise, we aggregate our firm-level estimates to assess the total effect of gold clause exposure on investment in 1933 and 1934. In this exercise, we abstract from general equilibrium effects and assume that, absent the uncertainty surrounding the abrogation, all firms would have invested as if they had no gold clause exposure (i.e., as if $\tilde{d}_j = 0$). We estimate the foregone aggregate investment rate in each year by computing the capital-weighted average gold clause effect:
\begin{equation}
\label{eq:aggregate}
\text{Gold clause effect}_t = \frac{\beta_t \sum_j \tilde{d}_{j,t} \times \text{Fixed capital}_{j,t-1}}{\sum_j \text{Fixed capital}_{j,t-1}},
\end{equation}
where $j$ indexes firms and $\beta_t$ is the coefficient on the interaction between $\tilde{d}_{j,t}$ and the indicator for year $t$.\footnote{To be included in this calculation, we require a firm to be present in the sample both in year $t$ and $t-1$.}

Table \ref{tab:agg} reports the aggregation results. Panel A presents estimates for all firms in our sample. The first row shows that the aggregate net investment rate was $-5.47\%$ in 1933 and $-2.47\%$ in 1934, indicating substantial divestment during the litigation period. The second row reports our baseline estimate of the gold clause effect: $-1.79$ and $-1.23$ percentage points in 1933 and 1934, respectively. These estimates imply that gold clause uncertainty accounts for roughly one-third of the aggregate divestment in 1933 ($1.79/5.47 = 33\%$) and approximately one-half in 1934 ($1.23/2.47 = 50\%$).

To provide context for these estimates, Panels B and C decompose total investment by gold clause exposure. Panel C shows that firms without gold-denominated debt ($\tilde{d}=0$) had aggregate net investment rates of $-4.10\%$ in 1933 and $-1.77\%$ in 1934. Panel B shows that firms with positive gold exposure ($\tilde{d}>0$) disinvested more severely, with aggregate net investment rates of $-6.59\%$ in 1933 and $-3.01\%$ in 1934, a gap of approximately $2.5$ and $1.2$ percentage points, respectively. 

The second row of Panel B refines this comparison by applying our regression-based estimates to gold-exposed firms only, yielding gold clause effects of $-3.27$ and $-2.18$ percentage points in 1933 and 1934, respectively---roughly half of these firms' total divestment in 1933 and approximately three-quarters in 1934. For both the full sample and the $\tilde{d}>0$ subsample, net investment improves substantially in the 1935--1940 period following the Supreme Court decision. The effect of gold exposure vanishes in later years, consistent with the view that the Supreme Court's decision resolved the surrounding leverage uncertainty.

If the investment response to gold exposure varies systematically across firm types---for example, by credit quality or size---then the baseline aggregation, which assumes a uniform marginal effect, could over- or underestimate the true aggregate impact. To address this concern, Internet Appendix Table \ref{tabapp:agg_het} reports results that allow for heterogeneous marginal effects by credit rating and firm size. Accounting for credit rating heterogeneity leaves the estimates largely unchanged. Allowing for size heterogeneity yields similar results in 1933 but a somewhat smaller aggregate effect in 1934, reflecting the concentration of capital among large firms that experienced more moderate gold clause effects. In both cases, the qualitative conclusions remain the same.

Finally, we interpret our aggregation results as providing a lower bound for the total effect of gold clause abrogation on private investment. As previously discussed, gold clauses were prevalent in other private contracts, including farm mortgages, public utility bonds, and railroad securities. While data limitations prevent us from directly studying the impact on these sectors, our results suggest that the effects may have been economically significant for these issuers as well.\footnote{To the best of our knowledge, disaggregated data on farm mortgage exposure to gold clauses are not readily available for this period. Additionally, most public utilities and railroads were not publicly traded during the 1930s, making it impossible to implement standard investment regressions for these firms.} Taken together, the partial-equilibrium aggregation results highlight the significant role of gold clause uncertainty in explaining the weak investment recovery of public firms in 1933 and 1934.
